% =========================================================
% Binary Operations
% =========================================================
\paragraph{Binary Operations}


% =========================================================
% Binary Operations --- Notes
% =========================================================

% ---------------------------------------------------------
% Definition --- Binary Operation
% ---------------------------------------------------------

\begin{definitionbox}{Definition (Binary Operation)}
\begin{definition}[Binary Operation]
\label{def:algebraic-group-binary-operation}
Let $G$ be a set. A \textbf{binary operation} on $G$ is a function
\[
\star : G \times G \to G.
\]
For $a, b \in G$, the element $\star(a,b)$ is denoted by $a \star b$.
\end{definition}
\end{definitionbox}

\begin{remark*}[Standard quantified statement]
\[
\begin{aligned}
&\star \text{ is a binary operation on } G \\
&\quad\Longleftrightarrow\quad
\forall a \in G\;\forall b \in G,\; a \star b \in G.
\end{aligned}
\]
\end{remark*}

\begin{remark*}[Definition predicate reading]
\[
\operatorname{BinaryOperation}(\star, G)
\;\colonequiv\;
\forall a \in G\;\forall b \in G,\; a \star b \in G.
\]
\end{remark*}

\begin{remark*}[Negated quantified statement]
\[
\begin{aligned}
&\star \text{ is not a binary operation on } G \\
&\quad\Longleftrightarrow\quad
\exists a \in G\;\exists b \in G \text{ such that } a \star b \notin G.
\end{aligned}
\]
\end{remark*}

\begin{remark*}[Negation predicate reading]
\[
\neg \operatorname{BinaryOperation}(\star, G)
\;\colonequiv\;
\exists a \in G\;\exists b \in G \text{ such that } a \star b \notin G.
\]
\end{remark*}

\begin{remark*}[Failure modes]
\[
\exists a \in G\;\exists b \in G \text{ such that } a \star b \notin G.
\]
\end{remark*}

\begin{remark*}[Failure mode decomposition]
\[
\begin{aligned}
&\star \text{ fails to be a binary operation on } G \\
&\quad\Longleftrightarrow\quad
\exists a \in G\;\exists b \in G \text{ such that } a \star b \notin G.
\end{aligned}
\]
\end{remark*}

\begin{remark*}[Interpretation]
A binary operation takes two elements of $G$ and returns another element of $G$.
This encodes closure of the operation within the set.
\end{remark*}

% ---------------------------------------------------------
% Definition --- Associativity
% ---------------------------------------------------------

\begin{definitionbox}{Definition (Associativity)}
\begin{definition}[Associativity]
\label{def:algebraic-group-associativity}
Let $\star$ be a binary operation on $G$. The operation $\star$ is
\textbf{associative} if
\[
(a \star b) \star c = a \star (b \star c)
\quad \text{for all } a, b, c \in G.
\]
\end{definition}
\end{definitionbox}

\begin{remark*}[Standard quantified statement]
\[
\begin{aligned}
&\star \text{ is associative on } G \\
&\quad\Longleftrightarrow\quad
\forall a \in G\;\forall b \in G\;\forall c \in G,\;
(a \star b) \star c = a \star (b \star c).
\end{aligned}
\]
\end{remark*}

\begin{remark*}[Definition predicate reading]
\[
\operatorname{Associative}(\star, G)
\;\colonequiv\;
\forall a \in G\;\forall b \in G\;\forall c \in G,\;
(a \star b) \star c = a \star (b \star c).
\]
\end{remark*}

\begin{remark*}[Negated quantified statement]
\[
\begin{aligned}
&\star \text{ is not associative on } G \\
&\quad\Longleftrightarrow\quad
\exists a \in G\;\exists b \in G\;\exists c \in G \text{ such that } \\
&\qquad (a \star b) \star c \ne a \star (b \star c).
\end{aligned}
\]
\end{remark*}

\begin{remark*}[Negation predicate reading]
\[
\neg \operatorname{Associative}(\star, G)
\;\colonequiv\;
\exists a \in G\;\exists b \in G\;\exists c \in G \text{ such that }
(a \star b) \star c \ne a \star (b \star c).
\]
\end{remark*}

\begin{remark*}[Failure modes]
\[
\exists a \in G\;\exists b \in G\;\exists c \in G \text{ such that }
(a \star b) \star c \ne a \star (b \star c).
\]
\end{remark*}

\begin{remark*}[Failure mode decomposition]
\[
\begin{aligned}
&\star \text{ fails associativity} \\
&\quad\Longleftrightarrow\quad
\exists a,b,c \in G \text{ witnessing a parenthesization mismatch.}
\end{aligned}
\]
\end{remark*}

\begin{remark*}[Interpretation]
Associativity means that the way elements are grouped does not affect the result.
\end{remark*}

% ---------------------------------------------------------
% Definition --- Commutativity
% ---------------------------------------------------------

\begin{definitionbox}{Definition (Commutativity)}
\begin{definition}[Commutativity]
\label{def:algebraic-group-commutativity}
Let $\star$ be a binary operation on $G$. The operation $\star$ is
\textbf{commutative} if
\[
a \star b = b \star a
\quad \text{for all } a, b \in G.
\]
\end{definition}
\end{definitionbox}

\begin{remark*}[Standard quantified statement]
\[
\begin{aligned}
&\star \text{ is commutative on } G \\
&\quad\Longleftrightarrow\quad
\forall a \in G\;\forall b \in G,\; a \star b = b \star a.
\end{aligned}
\]
\end{remark*}

\begin{remark*}[Definition predicate reading]
\[
\operatorname{Commutative}(\star, G)
\;\colonequiv\;
\forall a \in G\;\forall b \in G,\; a \star b = b \star a.
\]
\end{remark*}

\begin{remark*}[Negated quantified statement]
\[
\begin{aligned}
&\star \text{ is not commutative on } G \\
&\quad\Longleftrightarrow\quad
\exists a \in G\;\exists b \in G \text{ such that } a \star b \ne b \star a.
\end{aligned}
\]
\end{remark*}

\begin{remark*}[Negation predicate reading]
\[
\neg \operatorname{Commutative}(\star, G)
\;\colonequiv\;
\exists a \in G\;\exists b \in G \text{ such that } a \star b \ne b \star a.
\]
\end{remark*}

\begin{remark*}[Failure modes]
\[
\exists a \in G\;\exists b \in G \text{ such that } a \star b \ne b \star a.
\]
\end{remark*}

\begin{remark*}[Failure mode decomposition]
\[
\begin{aligned}
&\star \text{ fails commutativity} \\
&\quad\Longleftrightarrow\quad
\exists a,b \in G \text{ where swapping inputs changes the result.}
\end{aligned}
\]
\end{remark*}

\begin{remark*}[Interpretation]
Commutativity means the order of the inputs does not affect the result.
\end{remark*}

% ---------------------------------------------------------
% Remark --- Relationship
% ---------------------------------------------------------

\begin{remark*}[Interpretation]
Associativity and commutativity are independent properties. For example,
addition on $\mathbb{Z}$ is both associative and commutative; matrix
multiplication is associative but not commutative; subtraction on
$\mathbb{Z}$ is neither.
\end{remark*}































